\chapter{Tourette's}
\index{Tourette's}
\medskip
\noindent\textit{Educational reference; always seek professional advice for individual care.}

\section*{Definition and Overview}
Tourette's, commonly used shorthand for Tourette syndrome (also called Tourette's disorder), refers to a chronic tic disorder with both multiple motor tics and one or more vocal tics over a period exceeding one year, beginning in childhood or adolescence. In everyday local speech, families and schools often say ``Tourette's'' even when the formal label is still being clarified, or when a young person has chronic motor tics without prominent vocal tics. Precise terminology matters for research and services, but compassionate communication matters more in first consultations. People with Tourette's experience semi-voluntary movements and sounds that can be suppressed only briefly and usually rebound. The condition is neurological and developmental, not a reflection of manners, intelligence, or parenting. Public myths---especially that everyone with Tourette's swears uncontrollably---cause shame and delayed care. This chapter uses the familiar term Tourette's while describing accurate clinical features, supports, and pathways used in local health and education systems. Additional clinical context emphasises individualised care, shared decision-making, culturally safe communication, and alignment with current local guidelines. People should review written action plans after any change in medicines, and carers should know how to access after-hours advice. Documentation of goals, adverse effects, and follow-up timing improves continuity across general practice and specialist services. Educational material cannot replace personalised assessment when symptoms are severe, atypical, pregnancy-related, or occurring in infants, older adults, or immunocompromised people. Additional clinical context emphasises individualised care, shared decision-making, culturally safe communication, and alignment with current local guidelines. People should review written action plans after any change in medicines, and carers should know how to access after-hours advice. Documentation of goals, adverse effects, and follow-up timing improves continuity across general practice and specialist services. Educational material cannot replace personalised assessment when symptoms are severe, atypical, pregnancy-related, or occurring in infants, older adults, or immunocompromised people. Additional clinical context emphasises individualised care, shared decision-making, culturally safe communication, and alignment with current local guidelines. People should review written action plans after any change in medicines, and carers should know how to access after-hours advice. Documentation of goals, adverse effects, and follow-up timing improves continuity across general practice and specialist services. Educational material cannot replace personalised assessment when symptoms are severe, atypical, pregnancy-related, or occurring in infants, older adults, or immunocompromised people. Additional clinical context emphasises individualised care, shared decision-making, culturally safe communication, and alignment with current local guidelines. People should review written action plans after any change in medicines, and carers should know how to access after-hours advice. Documentation of goals, adverse effects, and follow-up timing improves continuity across general practice and specialist services. Educational material cannot replace personalised assessment when symptoms are severe, atypical, pregnancy-related, or occurring in infants, older adults, or immunocompromised people.

\section*{Epidemiology}
Using broad community screening, chronic tic disorders affect roughly one percent of children, with full Tourette's criteria met somewhat less often depending on study methods. Male predominance is consistent across countries. Onset before ten years is typical. Many children who have provisional tics never progress to chronic Tourette's; watchful waiting with support is appropriate when impairment is low. local data are incomplete because coding in primary care under-captures mild disease, yet clinical demand for behavioural therapy exceeds supply in several states. Indigenous and rural families may face longer waits for paediatric neurology and psychology. Co-occurrence with ADHD is among the highest of any neurodevelopmental pairing seen in community paediatrics. Transition to adult services is a recognised gap: adolescents who still have impairing tics can fall between paediatric and adult mental health or neurology services. Workplace disclosure rates remain low because of stigma, so adult prevalence appears lower than true rates.

\section*{Aetiology and Pathophysiology}
\begin{itemize}
    \item \textbf{Genetic contribution:} highly heritable polygenic risk with incomplete penetrance and variable expression within families
    \item \textbf{Basal ganglia networks:} impaired inhibitory control over motor patterns is central to current models
    \item \textbf{Neurotransmitter systems:} dopamine is most implicated; histamine, glutamate, and GABA pathways are also researched
    \item \textbf{Sensory-motor coupling:} premonitory urges show that tics are not purely random motor discharges
    \item \textbf{Developmental window:} symptoms emerge while motor control circuits are still maturing
    \item \textbf{Triggers versus causes:} sleep deprivation and stress amplify expression but do not create the underlying predisposition
\end{itemize}
No brain scan ``proves'' Tourette's. Pathophysiology research informs treatment targets but does not replace clinical diagnosis.

\section*{Clinical Features}
\begin{itemize}
    \item \textbf{Simple motor tics:} eye blinks, nose twitches, head flicks, shoulder jerks
    \item \textbf{Complex motor tics:} gestures, squatting, twirling, self-touching sequences, copropraxia in a minority
    \item \textbf{Simple vocal tics:} sniffs, grunts, coughs, throat clears, syllables
    \item \textbf{Complex vocal tics:} words, phrases, repeating heard speech, involuntary socially sensitive utterances in a minority
    \item \textbf{Urge--tic--relief cycle:} builds tension, discharges with the tic, briefly relieves, then returns
    \item \textbf{Context sensitivity:} may lessen during absorbing focused activity and worsen in evaluative social settings
    \item \textbf{Associated neurodevelopmental profile:} ADHD, OCD traits, sensory sensitivities, emotional dysregulation
\end{itemize}

\section*{Differential Diagnosis}
\begin{itemize}
    \item \textbf{Habitual behaviours:} nail biting or hair twirling without true premonitory tic urges
    \item \textbf{Allergic sniffing or chronic rhinitis:} medical review if purely nasal symptoms
    \item \textbf{Asthma cough:} nocturnal or exertional pattern differs from phonic tics
    \item \textbf{Functional neurological tics:} different onset story and variability
    \item \textbf{Stereotypic movement disorder:} especially with global developmental delay
    \item \textbf{Substance effects:} caffeine excess, stimulants, withdrawal states
    \item \textbf{Seizures:} usually distinguished by awareness changes and EEG when uncertain
\end{itemize}

\section*{When to Seek Medical Care}
Book a GP or paediatric review if movements or sounds are recurring, attracting negative attention, causing pain, or linked with attention, mood, or obsessive symptoms. Ask for referral to a clinician experienced in tic disorders when school function drops or when first-line advice is insufficient.

\textbf{Contact your local emergency services} for:
\begin{itemize}
    \item \textbf{Immediate safety concerns:} suicidal ideation with plan, severe aggression, or inability to be supervised safely
    \item \textbf{Trauma:} injuries from violent head or neck tics
    \item \textbf{Sudden severe neurological change:} weakness, unequal pupils, collapse---tics do not cause stroke-like syndromes, so new deficits need emergency assessment
    \item \textbf{Medication overdose:} accidental or intentional ingestion of tic or ADHD medicines
\end{itemize}

\section*{Diagnosis and Investigations}
\begin{itemize}
    \item \textbf{Clinical interview:} history from young person and caregivers separately when appropriate
    \item \textbf{Direct observation:} including video offered by families if tics are situational
    \item \textbf{Standard criteria:} duration, tic types, age at onset, exclusion of substances and other medical causes
    \item \textbf{Screening tools:} ADHD and OCD questionnaires as adjuncts, not standalone diagnoses
    \item \textbf{Investigations:} rarely needed for classic presentations
    \item \textbf{Educational assessment:} if learning difficulties coexist
\end{itemize}

\section*{Management}
\begin{itemize}
    \item \textbf{Watchful support:} for mild, non-impairing tics with education alone
    \item \textbf{Behavioural therapy:} CBIT where available; telehealth options expanding in many countries
    \item \textbf{Classroom strategies:} predictable routines, reduced call-outs for tics, safe exit passes
    \item \textbf{Parent training:} responding calmly, reinforcing competing responses practised in therapy
    \item \textbf{Medication:} when impairment is moderate--severe; start low, go slow, monitor
    \item \textbf{Peer support:} age-appropriate groups and credible online communities
    \item \textbf{Transition planning:} prepare adolescents for self-advocacy in TAFE, university, and work
\end{itemize}

\section*{Complications}
\begin{itemize}
    \item \textbf{Chronic pain:} from repetitive cervical tics
    \item \textbf{Social exclusion:} and cyberbullying
    \item \textbf{Anxiety amplification:} performance fear worsening tics in a vicious cycle
    \item \textbf{Academic disengagement:} especially if ADHD untreated
    \item \textbf{Self-esteem injury:} internalised stigma
    \item \textbf{Polypharmacy risk:} if multiple clinicians prescribe without coordination
\end{itemize}

\section*{Prognosis and Outcomes}
Most children with Tourette's become adults who work, form relationships, and participate fully in community life. A meaningful minority need ongoing psychological or pharmacological support. Outcomes improve when schools collaborate early and when comorbidities are treated assertively. Adults diagnosed late often describe relief at finally having a non-judgemental explanation for lifelong habits others criticised. Prognosis should be framed with cautious optimism rather than either catastrophe or dismissive minimisation. Additional clinical context emphasises individualised care, shared decision-making, culturally safe communication, and alignment with current local guidelines. People should review written action plans after any change in medicines, and carers should know how to access after-hours advice. Documentation of goals, adverse effects, and follow-up timing improves continuity across general practice and specialist services. Educational material cannot replace personalised assessment when symptoms are severe, atypical, pregnancy-related, or occurring in infants, older adults, or immunocompromised people. Additional clinical context emphasises individualised care, shared decision-making, culturally safe communication, and alignment with current local guidelines. People should review written action plans after any change in medicines, and carers should know how to access after-hours advice. Documentation of goals, adverse effects, and follow-up timing improves continuity across general practice and specialist services. Educational material cannot replace personalised assessment when symptoms are severe, atypical, pregnancy-related, or occurring in infants, older adults, or immunocompromised people. Additional clinical context emphasises individualised care, shared decision-making, culturally safe communication, and alignment with current local guidelines. People should review written action plans after any change in medicines, and carers should know how to access after-hours advice. Documentation of goals, adverse effects, and follow-up timing improves continuity across general practice and specialist services. Educational material cannot replace personalised assessment when symptoms are severe, atypical, pregnancy-related, or occurring in infants, older adults, or immunocompromised people. Additional clinical context emphasises individualised care, shared decision-making, culturally safe communication, and alignment with current local guidelines. People should review written action plans after any change in medicines, and carers should know how to access after-hours advice. Documentation of goals, adverse effects, and follow-up timing improves continuity across general practice and specialist services. Educational material cannot replace personalised assessment when symptoms are severe, atypical, pregnancy-related, or occurring in infants, older adults, or immunocompromised people.

\section*{Prevention and Self-Care}
\begin{itemize}
    \item \textbf{Protect sleep:} and reduce chaotic schedules where possible
    \item \textbf{Exercise regularly:} many people report short-term tic relief after sport
    \item \textbf{Plan high-stress events:} with extra rest beforehand
    \item \textbf{Use therapy homework:} daily short practice beats occasional long sessions
    \item \textbf{Coordinate clinicians:} one shared medication list
    \item \textbf{Challenge myths:} politely correct coworkers or classmates when safe to do so
    \item \textbf{Carer self-care:} parent stress management improves household climate
\end{itemize} Additional clinical context emphasises individualised care, shared decision-making, culturally safe communication, and alignment with current local guidelines. People should review written action plans after any change in medicines, and carers should know how to access after-hours advice. Documentation of goals, adverse effects, and follow-up timing improves continuity across general practice and specialist services. Educational material cannot replace personalised assessment when symptoms are severe, atypical, pregnancy-related, or occurring in infants, older adults, or immunocompromised people. Additional clinical context emphasises individualised care, shared decision-making, culturally safe communication, and alignment with current local guidelines. People should review written action plans after any change in medicines, and carers should know how to access after-hours advice. Documentation of goals, adverse effects, and follow-up timing improves continuity across general practice and specialist services. Educational material cannot replace personalised assessment when symptoms are severe, atypical, pregnancy-related, or occurring in infants, older adults, or immunocompromised people. Additional clinical context emphasises individualised care, shared decision-making, culturally safe communication, and alignment with current local guidelines. People should review written action plans after any change in medicines, and carers should know how to access after-hours advice. Documentation of goals, adverse effects, and follow-up timing improves continuity across general practice and specialist services. Educational material cannot replace personalised assessment when symptoms are severe, atypical, pregnancy-related, or occurring in infants, older adults, or immunocompromised people. Additional clinical context emphasises individualised care, shared decision-making, culturally safe communication, and alignment with current local guidelines. People should review written action plans after any change in medicines, and carers should know how to access after-hours advice. Documentation of goals, adverse effects, and follow-up timing improves continuity across general practice and specialist services. Educational material cannot replace personalised assessment when symptoms are severe, atypical, pregnancy-related, or occurring in infants, older adults, or immunocompromised people.

\section*{Sources and Further Reading}
Reputable sources for further reading include the following.

\begin{itemize}
    \item Healthdirect Australia. \textit{Tourette syndrome.} https://www.healthdirect.gov.au/
    \item Better Health Channel. \textit{Tourette syndrome.} https://www.betterhealth.vic.gov.au/
    \item NHS. \textit{Tourette's syndrome.} https://www.nhs.uk/
    \item Mayo Clinic. \textit{Tourette syndrome.} https://www.mayoclinic.org/
    \item MSD Manual. \textit{Tic disorders.} https://www.msdmanuals.com/
    \item MedlinePlus. \textit{Tourette syndrome.} https://medlineplus.gov/
\end{itemize}
